% !TEX root = main.tex
\section{Introduction}
Agentic AI systems increasingly generate both code and the checks meant to establish trust in it, such as unit tests~\citep{haque2026agentictests} and formal specifications~\citep{yang2024formalreasoning,bursuc2025benchmarkvericodingformallyverified}.
Deciding whether such a check is itself correct is the test-oracle problem~\citep{barr2015oracle}, now posed for checks that an AI wrote.
When the same system writes the code and its checks, the checks can be gamed: coding agents modify or delete failing tests instead of fixing the code~\citep{baker2025monitoring,zhong2025impossiblebench}.
Running a check cannot settle whether it is right, because a passing test shows only that the code and the test agree.
Benchmarks of AI-written tests accordingly grade them against a human-written reference: a generated test counts as reproducing an issue only if it fails before and passes after the developers' fix~\citep{mundler2024swtbench}.
Writing a formal specification is the formal analogue of writing the tests that SWE-bench uses to grade a patch~\citep{jimenez2024swebench}, with a stronger guarantee.
A property-based test checks a property on sampled inputs~\citep{claessen2000quickcheck}, whereas a theorem asserts it for every input.

In Lean~4, these artifacts are code-adjacent formal objects: runnable programs paired with unit tests, formal specifications (theorem statements about the program), and, sometimes, proofs of those theorems.
In classical autoformalization, the target theorem is supplied, so evaluation can often reduce to proof success.
In agentic verification, the system must decide \emph{what should be true} before it can prove anything.
The Lean kernel can check a proof, but it cannot tell us whether the proposed theorem is complete, semantically aligned with the intended program, or merely a weak statement that happens to typecheck~\citep{lahiri2024evaluating,yang2024formalreasoning}.
No kernel check or test run supplies this judgment of intent, so it falls to a human oracle~\citep{barr2015oracle}: autoformalization benchmarks rely on expert review of statement correctness~\citep{azerbayev2023proofnet}.
As agents produce more candidate artifacts than humans can inspect, specification-quality evaluation becomes a scalable-oversight problem~\citep{bowman2022scalable,zhou2025graders}.
LLM judges already stand in for the human oracle at scale.
Open theorem-prover pipelines use them to filter over a million autoformalized Lean statements~\citep{ying2024leanworkbook,lin2025goedelprover,wang2025kiminaprover}, and they are increasingly used to evaluate AI-generated code~\citep{Zhuo2024ICEScore,TongZhang2024CodeJudge,he2025llmjudgese}.
Yet the reliability of these judges on formal statements has received little study~\citep{chen2026reform}.
This paper asks when such a judge of AI-written formal specifications can be trusted.

The standard reliability check is per-judge human validation: collect labels, measure agreement between each judge and human ratings, and keep the judge if the correlation is high enough~\citep{Liu2023GEval,Zheng2023MTBench}.
The check is necessary because judge--human agreement varies substantially across judge models, datasets, and evaluated properties~\citep{Kocmi2024HumanVsLLM}.
Its cost, however, grows with every new judge, prompt, model update, and domain, and each one asks for another round of expert annotation.
% TODO(internal-review): a reviewer may ask why a new judge on the SAME domain cannot simply reuse the existing labels (no fresh ratings needed).
% The amortization argument is strongest for new domains / artifact distributions and for repeated judge selection on one fixed labeled set (adaptive overfitting; e.g., Dwork et al. 2015, "The reusable holdout").
% Consider stating explicitly which of these the protocol amortizes.

We instead collect human labels once and test each new judge automatically against falsifiable properties whose correct direction is known by construction.
A trustworthy judge should give the maximum score to an artifact compared with itself, and lower scores as semantic bugs accumulate or as theorem statements and unit tests are removed.
Its scores should also stay stable across repeated calls and semantics-preserving rewrites.
The one-time human panel measures how strongly these properties track human agreement; afterwards, a new judge configuration is screened by running the property suite, with no new labels.

On a 72-pair validation set with eight judge configurations, our variance-aware trust index \(\mathrm{TI}_{\mathrm{var}}\) is associated with judge-level human agreement (Spearman \(0.833\) on averaged labels for Codex judges), and the association is positive for Codex, Claude Opus, and Qwen3-Coder judges alike.
Human agreement alone can miss a bad judge: a text-similarity grader agrees with humans about as well as properly prompted judges, yet \(\mathrm{TI}_{\mathrm{var}}\) ranks it below all of them.
These are validation-set associations; a locked held-out test with independent raters remains future work.

Our \textbf{contributions} are:
\begin{enumerate}
  \item \textbf{An amortized validation protocol:} a one-time human panel plus four falsifiable properties lets any new (model, prompt, repeats, aggregation) judge configuration be screened without new labels (Section~\ref{sec:property-protocol}).
  \item \textbf{Validation-set evidence:} across three judge families, \(\mathrm{TI}_{\mathrm{var}}\) is associated with judge-level human agreement, and specification monotonicity is the most informative single property (Section~\ref{sec:validation-human-alignment}).
\end{enumerate}

